\قسمت{سرویس‌های کوچولو}

\زیرقسمت{تراکنش‌ها}

از شما خواسته شده است یک اپلیکیشن مالی طراحی کنید، این اپلیکیشن تراکنش‌\پانویس{transaction}های زیادی دارد. تیم شما یک تیم کوچک است، آیا از معماری میکروسرویس استفاده می‌کنید؟

\زیرقسمت{رای‌گیری توزیع شده}

از شما خواسته شده است یک سامانه رای‌گیری طراحی کنید. این سامانه رای‌گیری در همه شهرها مستقر می‌شود و در پایتخت یک سامانه مرکزی دارد که همه داده‌ها را جمع می‌کند. مدت رای‌گیری یک روز کامل است و احتمال خرابی در شبکه بین شهرها وجود دارد. در نظر داشته باشید هیچگاه نباید عملیات رای‌گیری در یک شهر متوقف شود.
طراحی شما به این صورت است که همه شهرها رای‌گیری را انجام می‌دهند و نتایج را به پایتخت به صورت لحظه‌ای ارسال می‌کند. در صورتی که خطا شبکه داشته باشیم نتایج ذخیره شده و بعد از رفع خطا ارسال می‌شوند.

در صورتی که نتایج نهایی اهمیت داشته باشند آیا می‌توان به نتایج سامانه شما اطمینان کرد؟ توضیح دهید. در صورتی که بخواهیم به صورت زنده هم نتایج را داشته باشیم آیا راهی وجود دارد که این نتایج همواره صحیح باشند؟ (راهنمایی: به تئوری \متن‌لاتین{CAP} دقت کنید.)

\نمره{۳}

% پاسخ پیشنهادی است و پیش از استفاده در تصحیح نیاز به بازبینی دارد.
\begin{پاسخ}

\متن‌سیاه{تراکنش‌ها.} خیر.
با تیم کوچک و تراکنش‌های زیاد، شکستن سامانه به میکروسرویس یعنی تبدیل تراکنش‌های محلی پایگاه داده به تراکنش‌های توزیع‌شده و الگوهایی مانند \متن‌لاتین{Saga}،
به‌علاوه‌ی بار عملیاتی استقرار و پایش چند سرویس. شروع با یک تک‌پارچه‌ی مرتب و جدا کردن بعدی سرویس‌ها منطقی‌تر است.

\متن‌سیاه{رای‌گیری.} این طراحی \متن‌لاتین{AP} است: در زمان خطای شبکه رای‌گیری ادامه پیدا می‌کند و داده‌ها بعدا می‌رسند.
به نتیجه‌ی \متن‌سیاه{نهایی} می‌توان اطمینان کرد، به شرط آنکه داده‌ای گم نشود و ادغام آن‌ها تکرارپذیر\پانویس{Idempotent} باشد؛ یعنی سازگاری نهایی\پانویس{Eventual Consistency}.
اما نتیجه‌ی \متن‌سیاه{زنده} در زمان افراز شبکه نمی‌تواند همواره درست باشد؛ تنها می‌توان آن را به صورت «موقت» و با ذکر شهرهای گزارش‌نشده نمایش داد.

\شروع{فقرات}
\فقره پاسخ منفی به قسمت اول به همراه دلیل: ۵۰ درصد نمره
\فقره تفکیک نتیجه‌ی نهایی از نتیجه‌ی زنده در قسمت دوم: ۵۰ درصد نمره
\پایان{فقرات}

\end{پاسخ}
